\section{Conclusion}
\label{sec:conclusion}

\subsection{Summary}

We have compared all 15 implemented redistricting algorithms across both
structure and search layers on five states (NC $k=14$, WI $k=8$, TX $k=38$,
FL $k=28$, NH $k=2$) and five metrics (EC, PP, D-seat, ACF(100), runtime).
The main findings:

\begin{enumerate}

  \item \textbf{ILP certifies what heuristics approximate.}
        For NH ($k = 2$), ILP achieves EC = 47 (certified optimal), 8.5\%
        below Simulated Annealing (51) and 23.4\% below METIS (58).
        The existence of the ILP baseline quantifies the optimality gap
        of every other structure algorithm for the first time.

  \item \textbf{The structure layer determines the PP ceiling.}
        CVD-Geographic achieves the highest mean PP of any algorithm
        across all states, outperforming even SMC-Percentile (the best
        search method for PP).
        Search algorithms improve PP by at most 8\% from the structure
        baseline; a poor PP structure choice cannot be fully compensated
        by search.
        When PP is the primary legal criterion, CVD-Geographic structure
        should be selected first, and the search algorithm selected second.

  \item \textbf{Parallel Tempering dominates EC for large $k$.}
        For TX ($k=38$) and FL ($k=28$), PT achieves lower final EC than
        Short-Burst by escaping basin boundaries inaccessible to local
        burst strategies.
        The crossover from Short-Burst to PT is at approximately $k \approx
        20$; for smaller $k$, Short-Burst is sufficient and faster on
        single-core hardware.

  \item \textbf{Adaptive Multi-scale is the fastest-mixing chain for TX.}
        AMS reduces ACF(100) to 0.31 for TX vs.\ 0.37 for fixed
        Multi-scale, a 16\% improvement, by front-loading coarse steps
        during warm-up.
        For $k < 15$, AMS and fixed Multi-scale are indistinguishable.

  \item \textbf{The full portfolio is covered by six questions.}
        Two structure questions (certified optimality? PP priority?) and
        four search questions (VRA? $k \geq 15$? calibrated distribution?
        compactness optimisation?) route any redistricting task to the
        correct structure + search combination.

\end{enumerate}

\subsection{The Three-Layer Compositor as a Portfolio System}

The three-layer compositor (\texttt{--structure}, \texttt{--weights-override},
\texttt{--search}) was designed to make algorithms composable.
The comparison in this paper confirms that composability is not merely
a software design property but a substantive algorithmic advantage:
the best EC results require both a good structure algorithm (SA) and a
good search algorithm (PT); the best PP results require both a good structure
algorithm (CVD-Geo) and a calibrated search algorithm (SMC-Percentile).
Neither layer alone achieves the best of both metrics simultaneously, and
no single combination dominates all metrics.

This composability is the defining property of the platform.
Prior redistricting tools typically implement one algorithm end-to-end
(GerryChain's ReCom, PlanScore's baseline, redist's SMC).
The three-layer compositor turns algorithm selection into a parameter choice,
enabling practitioners to express their redistricting requirements directly
in the CLI and obtain the appropriate algorithm portfolio automatically.

\subsection{Future Work}

\paragraph{Confirmed interactions.}
The structure--search interaction (Section~\ref{sec:decision}) is the primary
open empirical question: does SA + PT compound (both improving EC), or does
SA's baseline already reduce PT's marginal gain?
A controlled experiment varying structure with fixed PT search, and varying
search with fixed SA structure, would resolve this.

\paragraph{Additional structure algorithms.}
B.25 (in planning) will introduce a spectral bisection structure algorithm
based on the Fiedler vector of the graph Laplacian.
Spectral bisection is theoretically related to minimum-\ec{} bisection and
may bridge the gap between METIS and ILP for moderate-$N$ instances.

\paragraph{Multi-run confidence intervals.}
All results in this paper are single-run.
A follow-up study with $R = 20$ independent runs per algorithm--state pair
would provide confidence intervals and enable formal statistical comparison
of the close results (e.g., PT vs.\ AMS for ACF; CVD-Geo vs.\ SMC-Percentile
for PP).

\paragraph{G.15 as practitioner reference.}
This paper is intended as the definitive practitioner reference for the
\texttt{redist} algorithm portfolio.
As new algorithms are added to the B/G/H-series papers, they will be
incorporated into the combined Table~\ref{tab:combined} and the six-question
decision tree.
The framework's modular structure is designed to accommodate these additions
without restructuring the existing decision logic.

\subsection{Implementation Reference}

All 15 algorithms are implemented in the \texttt{redist} Rust workspace:

\begin{itemize}
  \item \texttt{crates/redist-core/src/structure/} --- structure algorithms
        (\texttt{metis.rs}, \texttt{simulated\_annealing.rs},
        \texttt{cvd\_graph.rs}, \texttt{cvd\_geographic.rs},
        \texttt{bfs\_growth.rs})
  \item \texttt{crates/redist-ilp/} --- ILP redistricting crate (B.24)
  \item \texttt{crates/redist-ensemble/src/} --- search algorithms
        (\texttt{shortburst.rs}, \texttt{flip.rs},
        \texttt{forest\_recom.rs}, \texttt{merge\_split.rs},
        \texttt{multiscale.rs}, \texttt{shortburst\_forest.rs},
        \texttt{vra\_recom.rs}, \texttt{parallel\_tempering.rs},
        \texttt{smc.rs}, \texttt{smc\_percentile.rs},
        \texttt{adaptive\_multiscale.rs})
  \item \texttt{crates/redist-cli/src/compositor/} --- three-layer compositor
        and CLI flag routing
\end{itemize}

The decision framework in Table~\ref{tab:combined} is the recommended
reference for algorithm selection in the \texttt{redist} platform.
All runs in this paper are reproducible from seed 42 and the configurations
in Section~\ref{sec:setup}.
